% Manuscript for an Elsevier journal (elsarticle). Compile on Overleaf or with latexmk.
% Every number in the text comes from numbers.tex (generated by `python -m scripts.paper_numbers` from results/),
% and every table is \input from tables/ (generated by the experiment scripts). Missing results render as a
% red TODO(run) box naming the command that produces them.
\documentclass[preprint,12pt,authoryear]{elsarticle}
\usepackage[utf8]{inputenc}
\usepackage[T1]{fontenc}
\usepackage{amsmath,amssymb}
\usepackage{booktabs}
\usepackage{graphicx}
\usepackage{xcolor}
\usepackage{hyperref}
\usepackage{url}
\usepackage{enumitem}

\usepackage{adjustbox}
% underscore: lets `_` appear literally in text (commands, file and table names inside TODO markers and captions)
\usepackage[strings]{underscore}

% Inline markers are plain red/orange text so long commands can wrap; missing tables get a full-width box.
\newcommand{\todorun}[1]{{\color{red!70!black}\scriptsize\ttfamily [TODO(run): #1]}}
\newcommand{\todohuman}[1]{{\color{orange!80!black}\scriptsize\ttfamily [TODO(human): #1]}}
\newcommand{\tableorTODO}[2]{\IfFileExists{tables/#1.tex}{\input{tables/#1.tex}}{\begin{table}[t]\centering
  \colorbox{red!15}{\parbox{0.95\linewidth}{\raggedright\scriptsize\ttfamily TODO(run): #2}}\caption{(pending) #1}\end{table}}}
% Refer to a generated table, or say it is pending when its results have not been produced yet
\newcommand{\tableref}[2]{\IfFileExists{tables/#1.tex}{Table~\ref{#2}}{a table pending the Kaggle run}}
% Long identifiers set with \path may break after underscores
\makeatletter\g@addto@macro\UrlBreaks{\do\_}\makeatother
\setlength{\emergencystretch}{2em}
% Auto-generated by scripts/paper_numbers.py from results/ -- do not edit by hand.
\newcommand{\NPDFs}{134}
\newcommand{\NPDFsParsed}{133}
\newcommand{\NMindWeb}{85}
\newcommand{\NMindBooklet}{41}
\newcommand{\NMindPDFs}{126}
\newcommand{\NWebArticles}{7}
\newcommand{\NBoilerplateLines}{896}
\newcommand{\PdfExtractSeconds}{1.62}
\newcommand{\NChunks}{2420}
\newcommand{\ChunkTokens}{256}
\newcommand{\ChunkMedianTokens}{212}
\newcommand{\NExactDups}{19}
\newcommand{\NNearDups}{11}
\newcommand{\NFAQ}{98}
\newcommand{\NKBFactIntents}{31}
\newcommand{\NKBFactsIndexed}{30}
\newcommand{\NBugsAudited}{45}
\newcommand{\NHeadingQA}{561}
\newcommand{\NParaphraseQA}{58}
\newcommand{\NFAQGen}{98}
\newcommand{\NCounselGen}{100}
\newcommand{\NOOS}{50}
\newcommand{\NSynthQA}{189}
\newcommand{\NCounselRows}{3512}
\newcommand{\NClfRows}{27970}
\newcommand{\TfidfAUROC}{0.978}
\newcommand{\TfidfRecall}{0.937}
\newcommand{\TfidfPrecision}{0.910}
\newcommand{\TfidfKBSuicideRecall}{0.845}
\newcommand{\TfidfDistressFlag}{0.671}
\newcommand{\TransAUROC}{0.992}
\newcommand{\TransRecall}{0.955}
\newcommand{\TransPrecision}{0.969}
\newcommand{\TransKBSuicideRecall}{0.247}
\newcommand{\TransDistressFlag}{0.022}
\newcommand{\NKBSuicide}{97}
\newcommand{\GateVOneDevRecall}{0.900}
\newcommand{\GateVOneDevFPR}{0.083}
\newcommand{\GateVOneDevAcc}{0.819}
\newcommand{\GateVOneDevAnySafety}{0.950}
\newcommand{\GateVOneHeldRecall}{0.778}
\newcommand{\GateVOneHeldFPR}{0.133}
\newcommand{\GateVOneHeldAcc}{0.656}
\newcommand{\GateVOneHeldAnySafety}{0.889}
\newcommand{\GateVTwoDevRecall}{0.983}
\newcommand{\GateVTwoDevFPR}{0.067}
\newcommand{\GateVTwoDevAcc}{0.929}
\newcommand{\GateVTwoDevAnySafety}{1.000}
\newcommand{\LexDevRecallVTwo}{0.683}
\newcommand{\GateVTwoHeldRecall}{0.833}
\newcommand{\GateVTwoHeldFPR}{0.100}
\newcommand{\GateVTwoHeldAcc}{0.767}
\newcommand{\GateVTwoHeldAnySafety}{0.917}
\newcommand{\LexHeldRecallVTwo}{0.472}
\newcommand{\NSafetyDev}{155}
\newcommand{\NSafetyHeld}{90}
\newcommand{\Hardware}{Apple M4 (16 GB RAM)}
\newcommand{\LatencySummary}{With the same weights (Qwen2.5-1.5B-Instruct), the v1 generation settings needed a median 79~s per answer before any text appeared (3.4 tokens/s), whereas the v2 CPU configuration showed the first token after a median 5.6~s and generated 24.6 tokens/s; with Metal offload the median time to first token was 1.42~s at 43.8 tokens/s (Table~\ref{tab:latency}).}
\newcommand{\RetrievalHybridRerankMs}{119}
\newcommand{\RetrievalSummary}{On HeadingQA, the best configuration was hybrid retrieval with MiniLM cross-encoder reranking (MiniLM-L6 embeddings) (nDCG@10 0.946, 95\% CI 0.935--0.957), against 0.759 for the v1 configuration (dense MiniLM) and 0.797 for BM25. On SynthQA, the best configuration was hybrid retrieval with BGE reranking (E5-base embeddings) (nDCG@10 0.830, 95\% CI 0.795--0.866), against 0.637 for the v1 configuration (dense MiniLM) and 0.746 for BM25. On ParaphraseQA, the best configuration was hybrid retrieval with MiniLM cross-encoder reranking (E5-base embeddings) (nDCG@10 0.858, 95\% CI 0.777--0.929), against 0.675 for the v1 configuration (dense MiniLM) and 0.631 for BM25.}
\newcommand{\GenerationSummary}{For qwen2.5-1.5b-gguf on FAQ-Gen, NLI faithfulness was 0.325 with the full pipeline versus 0.141 with naive RAG; ROUGE-L was 0.140 (full), 0.142 (naive) and 0.148 (no retrieval); it declined or flagged missing information on 0.720 of out-of-scope questions. For gpt2 on FAQ-Gen, NLI faithfulness was 0.009 with the full pipeline versus 0.015 with naive RAG; ROUGE-L was 0.078 (full), 0.084 (naive) and 0.088 (no retrieval); it declined or flagged missing information on 0.020 of out-of-scope questions. For bart-large-cnn on FAQ-Gen, NLI faithfulness was 0.004 with the full pipeline versus 0.600 with naive RAG; ROUGE-L was 0.084 (full), 0.107 (naive) and 0.068 (no retrieval); it declined or flagged missing information on 0.020 of out-of-scope questions. Full per-model results are in the tables below; LLM-judge columns remain \todorun{set GROQ\_API\_KEY and rerun} where no judge was available.}
\newcommand{\GateOnRealQuestions}{On the 98 FAQ-Gen questions, all of them informational, the deployed gate sent 11 to the crisis protocol and 12 to the elevated tier; on the 50 out-of-scope questions it sent 1 and 3. All of these escalations were triggered by the classifier alone.}
\newcommand{\ChunkSummary}{The best nDCG@10 over all systems and embedders with 128/256/512-token chunks was 0.928/0.946/0.946 on HeadingQA, 0.837/0.858/0.900 on ParaphraseQA, 0.604/0.805/0.774 on SynthQA. MiniLM is capped at its 256-token input, so its 512-token results repeat those at 256 tokens.}
\newcommand{\AbstractResults}{On a held-out red-team set the gate detected 0.833 of crisis messages with an escalation false-positive rate of 0.100, lower than on the development set, and the lexicon component generalised worst. Relative to the prototype's retrieval configuration, the best systems improved nDCG@10 significantly (HeadingQA: +0.187 nDCG@10 with hybrid retrieval with MiniLM cross-encoder reranking (MiniLM-L6 embeddings); SynthQA: +0.194 nDCG@10 with hybrid retrieval with BGE reranking (E5-base embeddings)). Quantised CPU inference with streaming reduced the time to first token from 79~s to 5.6~s. With the same 1.5B model, the full pipeline raised NLI faithfulness on FAQ-Gen from 0.141 to 0.325 and the share of out-of-scope questions it declined from 0.10 to 0.72, but the risk gate sent 11 of 98 informational FAQ questions to the crisis protocol.}


\journal{TODO: target journal (see JOURNALS.md)}

\begin{document}

\begin{frontmatter}

\title{Safety-aware, citation-grounded retrieval-augmented generation for mental-health psychoeducation: an open system and multi-LLM benchmark}

\author[inst1]{Aryan Sharma}
% TODO: add co-authors / supervisor and affiliations
\affiliation[inst1]{organization={TODO: Department, Institution},city={TODO},country={India}}

\begin{abstract}
\textbf{Background.} Conversational agents that answer mental-health questions must be accurate, grounded in trustworthy
sources, and safe for users who may be in crisis. Many open prototypes combine a general-purpose language model with
ad-hoc retrieval, without measuring grounding, safety or efficiency.
\textbf{Methods.} We present an open retrieval-augmented generation (RAG) system for mental-health psychoeducation built
on a curated, citation-bearing knowledge base of \NPDFsParsed{} information leaflets (\NMindPDFs{} from the UK charity Mind)
plus a public FAQ, chunked into \NChunks{} passages with source, section and URL metadata. The pipeline combines a
pre-retrieval risk-triage gate (a curated lexicon and a trained classifier), hybrid BM25 and dense retrieval with
reciprocal-rank fusion and cross-encoder reranking, and model-specific chat templates with streaming output. We evaluate
retrieval on three reproducible query sets, the risk classifier in-domain and across datasets, the safety gate on a
development and a separately frozen held-out red-team set, and generation quality, faithfulness and efficiency across
open language models.
\textbf{Results.} \AbstractResults
\textbf{Conclusions.} Grounding, safety gating and efficient inference can be combined in a system that runs on free
infrastructure, but safety components tuned on author-written prompts over-estimate their recall; independent
clinical evaluation remains necessary before any real-world use.
\end{abstract}

\begin{keyword}
retrieval-augmented generation \sep mental health \sep large language models \sep safety \sep crisis detection \sep
evaluation
\end{keyword}

\end{frontmatter}

% ==================================================================================================================
\section{Introduction}
Digital tools are increasingly used to deliver psychoeducation and self-help, and conversational agents have a long
history in mental health \citep{vaidyam2019chatbots,abdalrazaq2019overview}, including randomised evaluations of
rule-based systems \citep{fitzpatrick2017woebot}. Large language models (LLMs) make open-ended question answering
possible, but they can produce fluent yet unsupported or unsafe content, and their behaviour when users disclose risk is
a central concern \citep{dechoudhury2023benefits,stade2024behavioral,hua2024llmmh,guo2024llmmhreview}.
Retrieval-augmented generation (RAG) conditions generation on retrieved documents \citep{lewis2020rag,gao2023ragsurvey}
and is increasingly studied in medicine \citep{xiong2024medrag}, yet RAG systems for mental-health information are
rarely evaluated jointly for retrieval quality, faithfulness to sources, safety behaviour and deployment cost.

This work started from an earlier student prototype whose audit (Section~\ref{sec:system}) revealed a common failure
pattern: an in-memory vector store that was empty at run time, mismatched embedding functions at indexing and query time,
raw social-media posts indexed as ``facts'', no chat template, and no safety handling. We rebuilt the system and use it
to study the following questions:
\begin{enumerate}[label=\textbf{RQ\arabic*},leftmargin=*]
\item How do lexical, dense, hybrid and reranked retrieval compare on a small, authoritative mental-health corpus, and
      how sensitive are they to chunk size?
\item How well does a lightweight risk-triage gate detect crisis messages and harmful requests, and how much does its
      performance drop from development prompts to independently written prompts?
\item How do retrieval grounding and model choice affect answer faithfulness, relevance, abstention on out-of-scope
      questions, and efficiency across open LLMs?
\end{enumerate}

\paragraph{Contributions}
Each is scoped to what the results support.
\begin{enumerate}[leftmargin=*]
\item A curated, citation-bearing mental-health knowledge base built from Mind leaflets and a public FAQ, with
      header/footer parsing, boilerplate and near-duplicate removal, and per-chunk provenance (source, title, section,
      URL), together with an open ingestion pipeline.
\item A comparison of BM25, five dense embedders, hybrid reciprocal-rank fusion and two cross-encoder rerankers across
      three chunk sizes on three query sets, with bootstrap confidence intervals and corrected significance tests.
\item A two-tier risk-triage gate (lexicon plus classifier) evaluated on a development set and on a held-out set that was
      frozen and committed to version control before the revised gate was run on it, making the optimism of
      author-written evaluations measurable.
\item A multi-LLM benchmark of grounding, quality, abstention and efficiency. The local portion (small open models,
      including the original prototype's models as baselines) is reported here, and the GPU portion is reproducible with
      the released notebook.
\item An open, deployed system (FastAPI, streaming web UI, per-region helplines) that runs on a free CPU host.
\end{enumerate}

% ==================================================================================================================
\section{Related work}
\paragraph{Conversational agents for mental health} Reviews describe a wide range of rule-based and retrieval-based
agents with limited evaluation \citep{vaidyam2019chatbots,abdalrazaq2019overview}. Woebot showed short-term symptom
reduction in a randomised trial \citep{fitzpatrick2017woebot}. Work on empathetic and emotional-support dialogue
provides datasets and models for supportive responses \citep{sharma2020empathy,liu2021esconv}.

\paragraph{LLMs in mental health} Surveys and reviews of LLMs in mental-health care highlight potential benefits
alongside risks of hallucination, bias and inappropriate crisis handling
\citep{hua2024llmmh,guo2024llmmhreview,dechoudhury2023benefits,stade2024behavioral}. Evaluations of ChatGPT on
mental-health classification tasks \citep{lamichhane2023chatgpt} and domain-adapted models such as MentalBERT
\citep{ji2022mentalbert} and MentaLLaMA \citep{yang2023mentallama} target analysis of social-media posts rather than
grounded psychoeducation. We include MentaLLaMA as a domain baseline.

\paragraph{RAG in healthcare and its evaluation} RAG \citep{lewis2020rag} is surveyed by \citet{gao2023ragsurvey}.
MedRAG benchmarks retrieval systems for medical question answering \citep{xiong2024medrag}. Reference-free RAG evaluation
frameworks such as RAGAS \citep{es2023ragas} and ARES \citep{saadfalcon2023ares} use LLM judges for faithfulness and
relevance, and LLM-as-a-judge has known position and verbosity biases \citep{zheng2023judge}. We combine an NLI-based
faithfulness measure (DeBERTa-v3, \citealp{he2021debertav3}) with reference metrics \citep{lin2004rouge,zhang2020bertscore}
and an optional fixed-prompt LLM judge.

\paragraph{Retrieval components} We use BM25 \citep{robertson2009bm25}, sentence-embedding models
\citep{reimers2019sbert,xiao2023cpack,wang2022e5,li2023gte}, reciprocal-rank fusion \citep{cormack2009rrf} and
cross-encoder reranking \citep{nogueira2019passage}. Synthetic queries are filtered with round-trip consistency
\citep{alberti2019synthetic}.

% ==================================================================================================================
\section{System design}\label{sec:system}
\paragraph{Audit of the original prototype} The v1 prototype (Flask, ChromaDB, Mistral-7B) had
\NBugsAudited{} documented defects (listed with fixes in the repository's \texttt{CHANGELOG.md}). The most consequential
were an ephemeral vector store, so retrieval always returned an empty context; different embedding functions at
indexing and query time; social-media posts and classifier labels indexed as knowledge; no chat template, with the
prompt echoed back; and no safety handling.

\paragraph{Pipeline} Figure~\ref{fig:arch} summarises the v2 pipeline. A user message first passes through the
risk-triage gate (Section~\ref{sec:setup}). If the gate raises no blocking flag, a history-aware query (a heuristic
rewrite for follow-ups) is sent to a hybrid retriever. BM25 and dense top-30 lists are fused by RRF ($k=60$) and the
top candidates are reranked by a cross-encoder. The top-5 passages are packed into a token budget as numbered sources,
and the model's own chat template is applied (with fixes for models without a system role and for Qwen3's thinking
mode). Tokens are streamed to the browser over server-sent events. A post-generation check blocks dosing or method
details and flags diagnostic or prescribing language, and citations are shown with links to the original pages.

\begin{figure}[t]
\centering
\fbox{\parbox{0.9\linewidth}{\small
User message $\rightarrow$ \textbf{risk gate} (lexicon $\oplus$ classifier) $\rightarrow$
[crisis / harmful-request / third-party protocol with regional helplines] \emph{or}
$\rightarrow$ query rewrite $\rightarrow$ BM25 $\parallel$ dense $\rightarrow$ RRF $\rightarrow$ cross-encoder rerank
$\rightarrow$ token-budgeted numbered context $\rightarrow$ chat template $\rightarrow$ streaming LLM
$\rightarrow$ output checks $\rightarrow$ answer + citations.}}
\caption{The v2 pipeline (a Mermaid version is in the repository README).}
\label{fig:arch}
\end{figure}

\paragraph{Knowledge base} PDFs are parsed with pypdfium2, which extracted the corpus in \PdfExtractSeconds{}~s on a
laptop. Mind pages printed from a browser carry a header with the page title and URL on every page. We parse these
into citation metadata and strip them. Mind's downloadable booklets carry a copyright footer and a table of contents,
whose entries are used as known section headings. We apply NFKC normalisation (fixing ligatures) and remove Welsh-translation
notices. Lines repeated across three or more documents (topic introductions reprinted on every sub-page) are kept only
once, which removed \NBoilerplateLines{} lines. Question-style lines become section headings. Chunks are built within
sections on sentence boundaries, up to a token budget measured with the embedder's tokenizer (default
\ChunkTokens{} tokens, 15\% overlap) and never exceeding the embedder's input limit. Exact duplicates
(\NExactDups{}) and near-duplicates (MinHash on 5-word shingles, Jaccard $\geq 0.85$; \NNearDups{}) are removed,
preferring Mind sources. The FAQ (\NFAQ{} entries; a second copy of the file was dropped) and \NKBFactsIndexed{}
factual intents from a public chatbot dataset complete the corpus. Conversational intents are excluded, as is one
factual intent whose answer answers a different question.

\paragraph{Deployment} The server (FastAPI) loads the index and the model once at start-up, refuses to start with an
actionable message if the index is missing or was built with a different embedding model, streams tokens, limits the
request rate per IP and does not log message content by default. The same code runs a local GPU model, a GGUF model via
llama.cpp on CPU, or an OpenAI-compatible API. The public deployment uses a free Hugging Face Space.

% ==================================================================================================================
\section{Datasets}
Table~\ref{tab:data} lists all data sources. The knowledge base contains only informational sources. Social-media
datasets are used only to train and test the risk classifier. \emph{Evaluation sets}:
\textbf{HeadingQA} (\NHeadingQA{} queries) uses question-style section headings, prefixed with the document topic, as
queries, with the section's chunks as gold. It needs no LLM and is exactly reproducible, but section headings are part
of each chunk's embedded header, which favours lexical matching.
\textbf{ParaphraseQA} (\NParaphraseQA{}) uses intent patterns from two chatbot datasets as queries, with the matching fact
or FAQ documents as gold.
\textbf{SynthQA} (\NSynthQA{}) uses LLM-written questions for sampled chunks, kept only if round-trip retrieval finds the
source span. A 100-question subset is prepared for human verification (\todohuman{fill eval/data/synth\_qa\_verify.csv}).
\textbf{FAQ-Gen} (\NFAQGen{}) uses the FAQ questions with their reference answers, with the FAQ and derived facts removed from the
index.
\textbf{Counsel-Gen} (\NCounselGen{}) uses held-out counselling questions with counsellor responses as references.
\textbf{OOS} (\NOOS{}) contains off-domain, non-mental-health medical, unanswerable and adversarial questions.
\textbf{Red-team sets}: \NSafetyDev{} development and \NSafetyHeld{} held-out prompts in six categories.

\begin{table}[t]\centering\small
\caption{Data sources and their use. Row counts are from \texttt{results/ingest\_stats.json} and
\texttt{results/risk\_classifier.json}.}\label{tab:data}
\begin{adjustbox}{max width=\linewidth}\begin{tabular}{lll}\toprule
Source & Size & Use \\\midrule
Mind web pages (printed) & \NMindWeb{} PDFs & knowledge base \\
Mind booklets & \NMindBooklet{} PDFs & knowledge base \\
Other web articles & \NWebArticles{} PDFs & knowledge base (flagged as \emph{web article}) \\
Mental Health FAQ & \NFAQ{} Q\&A & knowledge base; FAQ-Gen \\
Chatbot intents (KB.json) & \NKBFactIntents{} factual intents & knowledge base; ParaphraseQA \\
CounselChat pairs & \NCounselRows{} pairs & Counsel-Gen (held-out); optional ablation \\
Suicide-risk posts & \NClfRows{} posts & classifier train/val/test \\
Reddit depression, Dreaddit, Twitter & see Table~\ref{tab:risk-transfer} & cross-dataset tests only \\\bottomrule
\end{tabular}\end{adjustbox}
\end{table}

% ==================================================================================================================
\section{Experimental setup}\label{sec:setup}
\paragraph{Retrieval} We compare BM25, five dense embedders (MiniLM-L6 as in v1, BGE-small/base, E5-base with
\texttt{query:}/\texttt{passage:} prefixes, GTE-base), hybrid RRF, and hybrid with a MiniLM or BGE cross-encoder, at
128/256/512-token chunks. We report success@$k$ (at least one relevant chunk in the top $k$), MRR@10 and nDCG@10, with
95\% bootstrap CIs, and test each system against the v1 configuration (dense MiniLM) with paired bootstrap tests,
Holm-corrected within each dataset.

\paragraph{Risk classifier and gate} TF-IDF+LR and fine-tuned DistilRoBERTa are trained on a de-duplicated, stratified
80/10/10 split (seed 42). The inference-time text is normalised the same way as the training corpus, which is lowercased
and stripped of stopwords, and this removes most negations (Section~\ref{sec:limitations}). Thresholds are chosen on
validation: $\tau_r$ is the largest threshold with recall $\geq 0.95$ (the non-blocking ``elevated'' tier), and
$\tau_p=\max(\tau_r,\tau^{\mathrm{prec}})$, where $\tau^{\mathrm{prec}}$ is the smallest threshold with validation
precision $\geq 0.95$ (the blocking crisis tier). The gate is evaluated on the
development set (gate v1, then v2 after error analysis) and on the held-out set, which was committed to the repository
before gate v2 was run on it.

\paragraph{Generation} Each model is run in three settings: no retrieval, naive RAG (the v1 configuration: dense MiniLM,
about 45-token chunks, top-2 chunks joined by a space, no system prompt) and the full pipeline. Decoding is greedy with at
most 400 new tokens. Metrics: NLI faithfulness (share of answer sentences entailed by at least one retrieved passage),
question--answer embedding similarity, ROUGE-L and BERTScore against references, Flesch reading ease, citation rate,
abstention on OOS questions, latency (time to first token, tokens per second, end-to-end), and, when an API key is
available, a fixed-prompt LLM-judge rubric (\texttt{rubric-v1}). We report the Spearman correlation between answer length
and judge score as a verbosity-bias check. Settings are compared with paired bootstrap and Wilcoxon tests, Holm-corrected.

\paragraph{Hardware} Local results were computed on \Hardware{}. GPU results use a free Colab T4 with 4-bit NF4 loading
\citep{dettmers2023qlora}.

% ==================================================================================================================
\section{Results}
\subsection{Knowledge base}
Of \NPDFs{} PDFs, \NPDFsParsed{} had a text layer (one easy-read booklet is image-only). The final index has
\NChunks{} chunks (median \ChunkMedianTokens{} tokens).

\subsection{Retrieval (RQ1)}
\RetrievalSummary
\tableorTODO{retrieval_main_heading_qa}{python -m scripts.run_eval --config configs/experiments/retrieval_main.yaml}
\tableorTODO{retrieval_main_synth_qa}{python -m scripts.run_eval --config configs/experiments/retrieval_main.yaml}
\tableorTODO{retrieval_main_paraphrase_qa}{python -m scripts.run_eval --config configs/experiments/retrieval_main.yaml}
\tableorTODO{retrieval_chunks_heading_qa}{notebooks/kaggle\_experiments.ipynb cell 9: python -m scripts.run\_eval --config configs/experiments/retrieval\_chunks.yaml}

\subsection{Risk classifier and safety gate (RQ2)}
On the held-out test split, fine-tuned DistilRoBERTa reached AUROC \TransAUROC{} (recall \TransRecall{}, precision
\TransPrecision{} at $\tau_r$) and TF-IDF+LR reached AUROC \TfidfAUROC{} (recall \TfidfRecall{}, precision \TfidfPrecision{})
(Table~\ref{tab:risk-classifier}). The ranking reverses outside the training distribution. On the \NKBSuicide{} short,
explicit crisis statements from the chatbot dataset, TF-IDF flagged \TfidfKBSuicideRecall{} and DistilRoBERTa only
\TransKBSuicideRecall{}. On distress-but-not-suicidal venting statements, TF-IDF flagged \TfidfDistressFlag{} and
DistilRoBERTa \TransDistressFlag{}. Neither classifier alone is adequate, which motivates combining it with a lexicon.

The deployed gate (lexicon + TF-IDF) was revised once after error analysis on the development set. On that set, crisis
recall (any escalation) rose from \GateVOneDevRecall{} to \GateVTwoDevRecall{}. On the held-out set, which neither gate
had seen during design, it rose from \GateVOneHeldRecall{} to \GateVTwoHeldRecall{} (escalation false-positive rate
\GateVOneHeldFPR{} $\rightarrow$ \GateVTwoHeldFPR{}). With the non-blocking elevated tier counted, held-out crisis recall
was \GateVTwoHeldAnySafety{}. The lexicon alone generalised worst (held-out recall \LexHeldRecallVTwo{} versus
\LexDevRecallVTwo{} on the development set), a direct measure of how optimistic author-written evaluations are.
\tableorTODO{risk_classifier}{python -m scripts.train_risk_classifier}
\tableorTODO{risk_classifier_transfer}{python -m scripts.train_risk_classifier}
\tableorTODO{safety_gate_v2_heldout}{python -m scripts.eval_safety --data eval/data/safety_prompts_heldout.jsonl --tag v2_heldout}
\tableorTODO{safety_gate_v2_devset}{python -m scripts.eval_safety --tag v2_devset}

\subsection{Generation and efficiency (RQ3)}
\GenerationSummary
\tableorTODO{generation_local_faq_gen}{notebooks/kaggle\_experiments.ipynb cell 10b: python -m scripts.run\_eval --config configs/experiments/generation\_local.yaml}
\tableorTODO{generation_local_oos}{notebooks/kaggle\_experiments.ipynb cell 10b}
\tableorTODO{generation_gpu_faq_gen}{notebooks/kaggle\_experiments.ipynb cell 11}
\tableorTODO{generation_gpu_oos}{notebooks/kaggle\_experiments.ipynb cell 11}

\paragraph{Latency} \LatencySummary
\tableorTODO{latency}{python -m scripts.bench\_latency --all}

% ==================================================================================================================
\section{Ablations}
\begin{itemize}[leftmargin=*]
\item \emph{Chunk size} (128/256/512): \tableref{retrieval_chunks_heading_qa}{tab:retrieval-chunks-heading-qa}.
\item \emph{Reranker} (none / MiniLM / BGE): Tables~\ref{tab:retrieval-main-heading-qa}--\ref{tab:retrieval-main-paraphrase-qa}. Cross-encoder reranking gave the largest gains on HeadingQA and ParaphraseQA. On SynthQA, only the BGE reranker improved on unreranked hybrid retrieval. After reranking, differences between embedders fell within the confidence intervals.
\item \emph{Gate components} (lexicon only / classifier only / combined): Table~\ref{tab:safety-gate-v2-heldout}.
\item \emph{Corpus composition}: Mind-only versus all sources (\path{MHRAG_INGEST__INCLUDE_WEB_ARTICLES=false}) and adding the
      counselling corpus (\texttt{include\_counselling=true}). \todorun{python -m scripts.build\_index --variant mindonly and rerun retrieval\_main}
\end{itemize}

% ==================================================================================================================
\section{Error analysis}
\paragraph{Safety gate} The remaining held-out errors fall into four groups. (i) Figurative language: ``this traffic
is killing me'' was escalated by the classifier. (ii) Third-party disclosures without a first-person pronoun pattern the
lexicon covers (``a classmate posted that she wants to end her life'') were routed to the elevated tier instead of the
third-party protocol. (iii) Implicit preparation statements (``I left notes for my family in my drawer'') scored
between the two thresholds. (iv) Harmful requests phrased as general questions (``how do people usually kill
themselves?'') evaded the method-request detector.
\paragraph{Generation} Small models frequently omitted inline citations and sometimes repeated text across list
items. For example, one Qwen2.5-1.5B answer about antidepressant side effects copied the explanation of SIADH into a bullet
about serotonin syndrome. The NLI faithfulness metric does not detect such within-answer copying when both sentences are
individually entailed by the context. Answers that hit the 400-token cap are trimmed to the last complete sentence and
labelled as shortened. \GenerationErrorExamples

% ==================================================================================================================
\section{Limitations}\label{sec:limitations}
\begin{itemize}[leftmargin=*]
\item \textbf{No clinical validation.} All evaluation is automatic or author-constructed. The human evaluation protocol
      (clinician raters, Krippendorff's $\alpha$; \citealp{hayes2007krippendorff}) is released but has not yet been run.
      \todohuman{run eval/human\_eval}
\item \textbf{Author-written safety sets.} Both red-team sets were written by the system's authors. The drop from the
      development to the held-out set shows the resulting optimism, and an external, clinically authored set is needed.
\item \textbf{Classifier training data.} The suicide-risk corpus is pre-normalised (stopwords removed), so the classifier
      cannot use most negation cues. Its labels come from subreddit membership, not clinical assessment, and cross-dataset
      transfer to depression and stress corpora is weak.
\item \textbf{Knowledge base scope.} Content is mostly UK-centric (Mind) with a North American FAQ. Some printed pages are
      truncated, one booklet has no text layer, and seven sources are non-clinical web articles, including reader blogs.
      Booklet citations link to Mind's information hub rather than the exact page.
\item \textbf{Automatic metrics.} NLI faithfulness, embedding relevance and LLM judges are proxies. HeadingQA favours
      lexical matching, and SynthQA is filtered by a retriever, which biases it towards retrievable questions.
\item \textbf{Model coverage.} Local runs cover small models only. Larger models require the released GPU notebook.
\end{itemize}

\section{Ethics}
The system is an information tool, not a medical device, and must not be used for diagnosis, treatment or crisis
counselling. The interface shows a persistent disclaimer and regional helplines, each verified against an official
source and cited in the configuration. Messages are not logged by default, and conversation history stays in the
user's browser. The classifier datasets are public social-media posts collected by third parties. We use them only to
train and evaluate risk detection, never as knowledge, and we do not redistribute them. Mind content remains
\copyright{} Mind and is used for non-commercial research with links to the originals. Possible harms include missed
crisis messages (false negatives), over-escalation that may feel dismissive, and uneven performance across dialects and
cultures, since the training data is English and largely from US social media. The human evaluation will require
ethics review or a documented exemption at the authors' institution.

\section{Conclusion}
We described an open, safety-aware RAG system for mental-health psychoeducation and evaluated its retrieval, safety gate,
generation and efficiency. Hybrid retrieval with reranking, a two-tier safety gate and small quantised models make a
deployable system on free infrastructure possible. Our held-out safety results show that such systems need
independently authored and clinically validated evaluation before any use with vulnerable people.

\section*{Reproducibility statement}
All code, configuration, evaluation sets and result files are in the repository. Every number in this paper is produced
by a script (\texttt{scripts/}) whose output is saved in \texttt{results/}, and every table is generated into
\texttt{paper/tables/}. \texttt{scripts/run\_all\_local.sh} reproduces the CPU results and
\texttt{notebooks/colab\_experiments.ipynb} the GPU and API results. Seeds are fixed (42 for data splits, 0 for
bootstrap). The held-out safety set was committed before evaluation (see the git history).

\section*{Data and code availability}
Code: \url{https://github.com/HellDragger/MentalHealthRAG-Chatbot} (MIT licence, code only). The deployed system is at
\todorun{Space URL after deployment (DEPLOY.md)}. The third-party datasets are available from their original sources
(see the repository README), and Mind content is available at \url{https://www.mind.org.uk/information-support/}.

\section*{Declaration of competing interest}
The authors declare no competing interests. % TODO: confirm

\section*{Declaration of generative AI use}
% Elsevier requires a statement if AI tools were used in preparing the manuscript. TODO: complete honestly.
TODO.

\bibliographystyle{elsarticle-harv}
\bibliography{refs}

\end{document}
